\scanpage{27}

Donc divisant $O_{X,x}$ par un idéal défini
par $\unsure{n|h}$ \unsure{élts}, on \deleted{\unsure{démontre}} trouve par
le th général \deleted{\unsure{sur}} (\unsure{puisque} :
$\underline O_{X,x}$ intègre \unsure{formellt} \unsure{équidim.} de
dim $\geq K+1=(n+h)+1$ ] un anneau
qui est connexe en dim $\geq\deleted{\unsure{(n+h)}}\inserted{K-n}$
$=\unsure{k|h}$, cqfd.

\uline{Cor. 1} au th. 2 Variante
avec \deleted{\unsure{connexité}} ``géométriquement \deleted{\unsure{int}}
unibranche'', $*$ et fibres géométriques.

\uline{Cor. 2} Variante avec $f$ propre
\unsure{mais|univ.} ouvert.

\uline{Cor. 3} $X\subset P^n_k$, \deleted{$X$ \unsure{de dim}}
\inserted{les composantes}
de $X$ de dim $\geq h+1$, et $X$ connexe
en dim $\geq h$. Alors pour $H$ système
de $\ell$ hyperplans $H_1$--$H_\ell$, \inserted{de même,} \deleted{\unsure{et}}
$X\cap H_1\cap\cdots\cap H_\ell=Y$
\unsure{satisfait les} conditions analogues avec
$h-\ell$ au lieu de $h$.

\margin{Généralisation au cas d'un morphisme
$X\longrightarrow P^n_k$.
Peut-on généraliser de façon analogue
le th. local ? Facile, \unsure{en}
utilisant le \unsure{th.} \unsure{EGA},
\unsure{connexité} \unsure{initial}.
}

\margin{26}
